\section{Experimental Results}
\label{sec:results}

Experimental data were collected using three machines:
\begin{itemize}
    \item Machine $M$: MacBook Pro with \textbf{M1} and \textbf{32~GB RAM}, built-in storage, operating system \textbf{macOS~26}, file system \texttt{APFS};
    % TODO: Check locked frequency
    \item Machine $A$: \textbf{Intel Core i7-8750H} (6 cores, 12 threads with HT, locked at 2.2~GHz), \textbf{16~GB DDR4 RAM}, \textbf{128~GB SATA M.2} SSD, operating system \textbf{Arch Linux}, kernel \texttt{6.17.1}, file system \texttt{btrfs};
    % TODO: Check locked frequency and PCIe lanes
    \item Machine $H$: \textbf{Intel Core i5-12450H} (4 E-cores and 4 P-Cores, 12 threads with Hyper Threading, E-Cores locked at 1.5~GHz and P-Cores at 2~GHz), \textbf{16~GB DDR4 RAM}, \textbf{512~GB NVMe 8xPCIe3} SSD, operating system \textbf{Arch Linux}, kernel \texttt{6.17.1}, file system \texttt{xfs}.
\end{itemize}

The machines were connected via a gigabit Ethernet network with a latency of about 0{,}3~ms while under stress.

A purpose made program \texttt{ktb} was developed to stress the system using multiple connections.
We first recorded baseline values, by making each machine stress itself in both get and set operations
on a table with disabled throttling (simply set to a high enough number) and value sizes of 4~kB\@.

Table~\ref{tab:results-stress} shows the results and allows us to rank the machine as follows: $H > M \approx A$.

% TODO: Bandwidth o altra parola?
The scalability and replication of the system were then tested.
The following tests will all run \texttt{ktb} on machine H unless stated otherwise and will take in consideration only set operations.
Get operations are not considered because as the OS caches file reads, their results not correspond to true random access disk performance,
nonetheless measured req/s for get operations are reported in the table~\ref{tab:results-stress}.
We configured the machines for either a uniform key distribution, where each machine's bandwidth parameter was set to the same sufficiently high value,
or with a balanced key distribution, where each machine's bandwidth was still set sufficiently high but proportionally to the self stress results: $M$ and $A$ with the same value and $H$ with double that.

To test the key distribution and scalability of the system, we used a table configured
with enough bandwidth to saturate the test case, with parameters $N=3$, $R=1$, and $W=1$.
This configuration spreads requests across peers according to the machines' bandwidth proportions.
We expect the system to provide, in the uniform case, about 3 times the rate of the slowest machine, and in the balanced case about the sum of the machines rate.
% TODO Spiegare meglio i conti?
Testing shows that the system is able to scale and managed to get to 92\% and 83\% of the expected values.
The losses can in part be accounted for by considering the CPU overhead of running the stress itself
and the probabilistic nature of key distribution.

%N=1000, p_i = 1/3
%var = 1000 * 1/3 * 2/3 = 2/9 * 1000
%stdev ~ 15
%1 * stdev ~ 68% => 1 - 333/333+15 = 0,04310344828
%0,9166666667

%N=1000, p_1 = p_2 = 1/4, p_3 = 1/2
%var_1,2 = 1000 * 1/4 * 3/4 = 3/16 * 1000
%stdev_1,2 ~ 14
%1 * stdev ~ 68% => 1 - 250/250+14 = 0,05303030303
%var_3 = 1000 * 1/2 * 1/2 = 1/4 * 1000
%stdev_3 ~ 16
%1 * stdev ~ 68% => 1 - 500/500+16 = 0,03100775194
%0,8252427184

To test the replication of the system, a table created with enough bandwidth to fill the test case, $N=2$, $R=2$ and $W=2$
was used as it allows for reproducible performance ($N=R=W$) and it was considered more interesting than a simpler $N=3$, $R=3$ and $W=3$ case.
We expect the system to provide, in the uniform case, about $\frac32$ times the slowest machine,
and in the balanced case about the sum of the two slowest machines.
% TODO Spiegare meglio i conti?
Testing shows that the system is able to scale and managed to get to 100\% and 96\% of the expected values.
The remaining 5\% cannot be accounted for solely by considering the CPU overhead of running the stress itself,
but is due to the probabilistic nature of key distribution.

%N=1000, p_1,2,3 = 2/3
%var_1,2,3 = 1000 * 2/3 * 1/3 = 2/9 * 1000
%stdev ~ 15
%1 * stdev ~ 68% => 1 - 666/666+15 = 0,02202643172
%1

%N=1000, p_1,2 = 1/2, p_3=1
%var_1,2 = 1000 * 1/2 * 1/2 = 1/4 * 1000
%stdev ~ 16
%1 * stdev ~ 68% => 1 - 500/500+16 = 0,031
%var_3 = 0
%0,9591836735


\begin{table}
    \centering
    \caption{TODO}
    \label{tab:results-stress}
    \begin{tabular}{lrrr}
        \toprule
        Test case               & get (512~B) & get (4~kB) & set (4~kB) \\
        \midrule
        Machine $M$ self stress &             & 64.0k      & 5.0k       \\
        Machine $H$ self stress &             & 83.3k      & 10.8k      \\
        Machine $A$ self stress &             & 56.6k      & 4.8k       \\
        Uniform on (3, 1, 1)    & 102.0k      & 44.1k      & 13.2k      \\
        Balanced on (3, 1, 1)   & 97.0k       & 56.5k      & 17.0k      \\
        Uniform on (2, 2, 2)    & 67.3k       & 21.8k      & 7.2k       \\
        Balanced on (2, 2, 2)   & 62.4k       & 28.0k      & 9.4k       \\
        \bottomrule
    \end{tabular}
\end{table}

% TODO
%\footnote[1]{
%    Limited by Gigabit speeds.
%}

The last test focused on the real-time properties of the system, simulating a concrete use case.
With a balanced bandwidth configuration, a table was created with a bandwidth of 200 req/s and $N=3$, $R=3$, $W=3$.
To represent external interference, another table was created with sufficiently high bandwidth, $N=2$, $R=2$, and $W=2$.
The system was successfully able to isolate the target table from the outside noise, guaranteeing the allocated 200 req/s.
We also measured latencies when the target table is receiving less than 200 req/s with and without stress, as reported in table~\ref{tab:results-isolation}.
Overall, there is a slight increase in latencies, which is to be expected as no explicit latency control is implemented.

\begin{table}
    \centering
    \caption{TODO}
    \label{tab:results-isolation}
    \begin{tabular}{lrrrrrr}
        \toprule
        & Average ms & $P_{50}$ms & $P_{90}$ms & $P_{99}$ms & $P_{99.9}$ms \\
        \midrule
        Baseline on H        & 1.28       & 1.22       & 1.38       & 1.63       & 19.50        \\
        Stressed on H        & 2.35       & 2.14       & 2.78       & 6.12       & 37.87        \\
        Distributed Baseline & 4.18       & 3.22       & 6.75       & 16.09      & 37.55        \\
        Distributed Stressed & 12.10      & 11.85      & 19.37      & 48.29      & 92.33        \\
        \bottomrule
    \end{tabular}
\end{table}
